\documentclass[12pt]{article}
\usepackage{amsmath, amssymb, geometry, graphicx, tcolorbox, enumerate}
\geometry{a4paper, margin=1in}

\title{\textbf{Homework 15.3: Double Integrals in Polar Coordinates}}
\author{Tashfeen Omran}
\date{\today}

\begin{document}

\maketitle
\tableofcontents
\newpage

\section{Part 1: Comprehensive Introduction, Context, and Prerequisites}

\subsection{Conceptual Context and Motivation}
Historically, mathematicians in the 17th and 18th centuries like Newton, Leibniz, and Euler developed calculus to solve problems involving areas, volumes, and physical quantities (like mass and centers of gravity). While the Cartesian coordinate system ($x,y,z$) is excellent for rectangular regions, mathematicians quickly realized that it becomes highly complex and tedious when calculating quantities for circular, cylindrical, or radially symmetric regions. 

To solve this, the \textbf{polar coordinate system} was formalized. Instead of navigating a grid by moving left/right ($x$) and up/down ($y$), polar coordinates identify a point by its distance from the origin ($r$) and its angle from the positive x-axis ($\theta$). 

In multivariable calculus, evaluating a double integral over a circular region using Cartesian coordinates involves messy square roots in the integration limits (e.g., $y = \sqrt{a^2 - x^2}$). By transitioning to polar coordinates, the boundary of a circle simply becomes a constant ($r = a$), transforming complex integrals into straightforward ones with constant bounds.

\subsection{Core Concepts and Intuition}
A double integral $\iint_D f(x,y) \, dA$ calculates the volume under a surface $f(x,y)$ over a region $D$ in the $xy$-plane. 
When we change from Cartesian to polar coordinates, we use the standard transformations:
\[ x = r\cos(\theta), \quad y = r\sin(\theta), \quad r^2 = x^2 + y^2 \]
The most critical part of this transformation is understanding what happens to the area differential, $dA$. In rectangular coordinates, an infinitesimally small piece of area is a rectangle: $dA = dx \, dy$. 

In polar coordinates, an infinitesimally small piece of area is a "polar rectangle" bounded by two radii (separated by a tiny angle $d\theta$) and two circular arcs (separated by a tiny radius $dr$). The length of the inner arc is $r \, d\theta$, and its width is $dr$. Therefore, the area of this tiny polar patch is approximately a rectangle with dimensions $r \, d\theta$ and $dr$. This gives us the fundamental relation:
\[ dA = r \, dr \, d\theta \]
\textit{Warning:} The most common mistake in this chapter is forgetting to include that extra $r$!

\subsection{Key Formulas}
\begin{tcolorbox}
\textbf{1. Cartesian to Polar Conversions:}
\[ x = r\cos(\theta), \quad y = r\sin(\theta), \quad x^2 + y^2 = r^2 \]
\textbf{2. The Polar Double Integral Formula:}
\[ \iint_D f(x,y) \, dA = \int_{\alpha}^{\beta} \int_{h_1(\theta)}^{h_2(\theta)} f(r\cos(\theta), r\sin(\theta)) \, r \, dr \, d\theta \]
\end{tcolorbox}

\subsection{Prerequisites}
To master this chapter, you must be proficient in:
\begin{itemize}
    \item \textbf{Trigonometric Identities:} Specifically the half-angle formulas: $\sin^2(\theta) = \frac{1 - \cos(2\theta)}{2}$ and $\cos^2(\theta) = \frac{1 + \cos(2\theta)}{2}$.
    \item \textbf{U-Substitution:} You will frequently need to integrate expressions like $r e^{-r^2}$ or $r\cos(r^2)$, which require letting $u = r^2$.
    \item \textbf{Basic Graphing in Polar:} Knowing the shapes of $r = a$ (circle), $r = a\sin(\theta)$ (shifted circle), and basic sectors.
\end{itemize}

\newpage
\section{Part 2: Detailed Homework Solutions}

\subsection*{Problem 1}
\textbf{Question:} Decide whether to use polar or rectangular coordinates for the region $R$ shown, and write the iterated integral. The shaded region is bounded by a circular arc of radius 1 and spans from the line $y = -x$ in the fourth quadrant to the line $y = -x$ in the second quadrant.

\textbf{Solution:}
The region $R$ is a semi-circle of radius 1. The line $y = -x$ corresponds to an angle of $-\pi/4$ in the 4th quadrant and $3\pi/4$ in the 2nd quadrant. Because the region is circular and based at the origin, polar coordinates are vastly superior. 
The bounds are:
\[ 0 \le r \le 1 \]
\[ -\pi/4 \le \theta \le 3\pi/4 \]
Converting the integral:
\[ \iint_R f(x,y) \, dA = \int_{-\pi/4}^{3\pi/4} \int_0^1 f(r\cos(\theta), r\sin(\theta)) \, r \, dr \, d\theta \]
\textbf{Final Answer:} The 4th option: $\int_{-\pi/4}^{3\pi/4} \int_0^1 f(r\cos(\theta), r\sin(\theta)) \, r \, dr \, d\theta$

\subsection*{Problem 2}
\textbf{Question:} Decide whether to use polar or rectangular coordinates. The region $R$ is identical in shape to Problem 1 but has a radius of 3.

\textbf{Solution:}
Just as in Problem 1, this is a semi-circle bounded by $y = -x$. The radius extends to 3.
The bounds are:
\[ 0 \le r \le 3 \]
\[ -\pi/4 \le \theta \le 3\pi/4 \]
\[ \iint_R f(x,y) \, dA = \int_{-\pi/4}^{3\pi/4} \int_0^3 f(r\cos(\theta), r\sin(\theta)) \, r \, dr \, d\theta \]
\textbf{Final Answer:} The 5th option: $\int_{-\pi/4}^{3\pi/4} \int_0^3 f(r\cos(\theta), r\sin(\theta)) \, r \, dr \, d\theta$

\subsection*{Problem 3}
\textbf{Question:} Decide whether to use polar or rectangular coordinates. The region $R$ is an annulus (ring) between a circle of radius 8 and a circle of radius 10.

\textbf{Solution:}
Because the region is formed by two full circles centered at the origin, polar coordinates are ideal.
The angle bounds cover the entire circle: $0 \le \theta \le 2\pi$.
The radius bounds span from the inner edge to the outer edge: $8 \le r \le 10$.
\[ \iint_R f(x,y) \, dA = \int_0^{2\pi} \int_8^{10} f(r\cos(\theta), r\sin(\theta)) \, r \, dr \, d\theta \]
\textbf{Final Answer:} The 5th option: $\int_0^{2\pi} \int_8^{10} f(r\cos(\theta), r\sin(\theta)) \, r \, dr \, d\theta$

\subsection*{Problem 4}
\textbf{Question:} Sketch the region and evaluate the integral: $\int_{\pi/2}^{\pi} \int_0^{8\sin(\theta)} r \, dr \, d\theta$

\textbf{Solution:}
\textbf{Step 1: Identify the region.}
The boundary $r = 8\sin(\theta)$ is a circle of radius 4 centered at $(0, 4)$ on the y-axis. The limits of integration for $\theta$ are from $\pi/2$ to $\pi$, which restricts us to the second quadrant (the left half of this circle).
\textbf{Step 2: Evaluate the inner integral.}
\[ \int_0^{8\sin(\theta)} r \, dr = \left[ \frac{r^2}{2} \right]_0^{8\sin(\theta)} = \frac{64\sin^2(\theta)}{2} = 32\sin^2(\theta) \]
\textbf{Step 3: Evaluate the outer integral.}
\[ \int_{\pi/2}^{\pi} 32\sin^2(\theta) \, d\theta \]
Use the half-angle identity $\sin^2(\theta) = \frac{1 - \cos(2\theta)}{2}$:
\[ 32 \int_{\pi/2}^{\pi} \frac{1 - \cos(2\theta)}{2} \, d\theta = 16 \int_{\pi/2}^{\pi} (1 - \cos(2\theta)) \, d\theta \]
\[ = 16 \left[ \theta - \frac{\sin(2\theta)}{2} \right]_{\pi/2}^{\pi} \]
Evaluate at upper limit $\pi$: $16(\pi - 0) = 16\pi$
Evaluate at lower limit $\pi/2$: $16(\pi/2 - 0) = 8\pi$
\[ 16\pi - 8\pi = 8\pi \]
\textbf{Final Answer:} $8\pi$

\subsection*{Problem 5}
\textbf{Question:} Evaluate $\iint_D 5x^2y \, dA$ where $D$ is the top half of the disk with center the origin and radius 4.

\textbf{Solution:}
\textbf{Step 1: Convert region $D$ to polar bounds.}
Top half of a circle of radius 4 means $0 \le r \le 4$ and $0 \le \theta \le \pi$.
\textbf{Step 2: Convert the integrand.}
\[ 5x^2y = 5(r\cos(\theta))^2(r\sin(\theta)) = 5r^3\cos^2(\theta)\sin(\theta) \]
Don't forget the $dA = r \, dr \, d\theta$.
The full integral becomes:
\[ \int_0^{\pi} \int_0^4 \left( 5r^3\cos^2(\theta)\sin(\theta) \right) r \, dr \, d\theta = \int_0^{\pi} \int_0^4 5r^4\cos^2(\theta)\sin(\theta) \, dr \, d\theta \]
\textbf{Step 3: Evaluate the integral.}
Since the limits are constant and the function is separable, we can split it:
\[ \left( \int_0^{\pi} \cos^2(\theta)\sin(\theta) \, d\theta \right) \times \left( \int_0^4 5r^4 \, dr \right) \]
For the $r$-integral:
\[ \int_0^4 5r^4 \, dr = \left[ r^5 \right]_0^4 = 4^5 - 0 = 1024 \]
For the $\theta$-integral, use $u$-substitution with $u = \cos(\theta)$, $du = -\sin(\theta) \, d\theta$.
The bounds change from $\theta = 0 \to u = 1$ and $\theta = \pi \to u = -1$.
\[ \int_1^{-1} -u^2 \, du = \int_{-1}^1 u^2 \, du = \left[ \frac{u^3}{3} \right]_{-1}^1 = \frac{1}{3} - \left(-\frac{1}{3}\right) = \frac{2}{3} \]
Multiply the results together:
\[ 1024 \times \frac{2}{3} = \frac{2048}{3} \]
\textbf{Final Answer:} $\frac{2048}{3}$

\subsection*{Problem 6}
\textbf{Question:} Find the volume of the solid bounded by $z=0$ and the paraboloid $z = 4 - x^2 - y^2$. Fill in the blanks.

\textbf{Solution:}
The region $D$ in the $xy$-plane is found by setting $z=0$: $0 = 4 - x^2 - y^2 \implies x^2 + y^2 = 4$. This is a circle of radius 2.
\[ V = \int_0^{2\pi} \int_0^2 (4 - r^2)r \, dr \, d\theta = \int_0^{2\pi} d\theta \int_0^2 (4r - r^3) \, dr \]
Evaluate the inner integral first.
\[ \int (4r - r^3) \, dr = 2r^2 - \frac{r^4}{4} \]
So the first blank asks for the antiderivative expression evaluated from $0$ to $2$.
Evaluating from 0 to 2 gives:
\[ 2(2)^2 - \frac{2^4}{4} = 8 - \frac{16}{4} = 8 - 4 = 4 \]
The final volume:
\[ 2\pi \times 4 = 8\pi \]
\textbf{Final Answers for Blanks:}
Box 1 (Antiderivative inside bracket): $2r^2 - \frac{r^4}{4}$
Box 2 (Evaluated number inside bracket): $4$
Box 3 (Final volume): $8\pi$

\subsection*{Problem 7}
\textbf{Question:} Change $\iint_D 4x^2y \, dA$ to polar coordinates, where $D$ is the top half of a disk of radius 5, then evaluate. Fill in limits $A, B$ and the integrand.

\textbf{Solution:}
\textbf{Step 1: Identify limits and integrand.}
Top half of the disk means $0 \le \theta \le \pi$, so lower limit $A = 0$ and upper limit $B = \pi$.
The radius goes from $0$ to $5$.
The integrand is $4(r\cos(\theta))^2(r\sin(\theta)) \cdot r = 4r^4\cos^2(\theta)\sin(\theta)$.
\textbf{Step 2: Evaluate the integral.}
\[ \int_0^{\pi} \int_0^5 4r^4\cos^2(\theta)\sin(\theta) \, dr \, d\theta \]
Separate the integrals:
\[ \left( \int_0^{\pi} \cos^2(\theta)\sin(\theta) \, d\theta \right) \left( \int_0^5 4r^4 \, dr \right) \]
From Problem 5, we know $\int_0^{\pi} \cos^2(\theta)\sin(\theta) \, d\theta = \frac{2}{3}$.
For the $r$ integral:
\[ \int_0^5 4r^4 \, dr = \left[ \frac{4}{5}r^5 \right]_0^5 = \frac{4}{5}(3125) = 2500 \]
Multiply them:
\[ \frac{2}{3} \times 2500 = \frac{5000}{3} \]
\textbf{Final Answers:}
Integral box (integrand): $4r^4\cos^2(\theta)\sin(\theta)$
$A = 0$
$B = \pi$
Final Evaluation box: $\frac{5000}{3}$

\subsection*{Problem 8 \& 9}
\textit{Note: Problem 8 and Problem 9 feature the exact same integral and region, just formatted differently in WebAssign. The solution is identical for both.}

\textbf{Question:} Evaluate $\iint_R (4x - y) \, dA$, where $R$ is the region in the first quadrant enclosed by the circle $x^2 + y^2 = 4$ and the lines $x = 0$ and $y = x$.

\textbf{Solution:}
\textbf{Step 1: Determine bounds.}
Radius is bounded by the circle: $0 \le r \le 2$.
The line $y = x$ corresponds to $\theta = \pi/4$.
The line $x = 0$ (the positive y-axis) corresponds to $\theta = \pi/2$.
So, $\pi/4 \le \theta \le \pi/2$. 
For Problem 8, this means $A = \pi/4$ and $B = \pi/2$.
\textbf{Step 2: Set up the integrand.}
\[ (4x - y) \, dA = (4r\cos(\theta) - r\sin(\theta)) r \, dr \, d\theta = (4r^2\cos(\theta) - r^2\sin(\theta)) \, dr \, d\theta \]
\textbf{Step 3: Evaluate the integral.}
\[ \int_{\pi/4}^{\pi/2} \int_0^2 r^2(4\cos(\theta) - \sin(\theta)) \, dr \, d\theta \]
Separate the variables:
\[ \left( \int_{\pi/4}^{\pi/2} (4\cos(\theta) - \sin(\theta)) \, d\theta \right) \left( \int_0^2 r^2 \, dr \right) \]
Evaluate the $\theta$-integral:
\[ \left[ 4\sin(\theta) + \cos(\theta) \right]_{\pi/4}^{\pi/2} \]
Upper limit ($\pi/2$): $4(1) + 0 = 4$
Lower limit ($\pi/4$): $4\left(\frac{\sqrt{2}}{2}\right) + \frac{\sqrt{2}}{2} = 2\sqrt{2} + \frac{\sqrt{2}}{2} = \frac{5\sqrt{2}}{2}$
$\theta$-integral result: $4 - \frac{5\sqrt{2}}{2}$
Evaluate the $r$-integral:
\[ \left[ \frac{r^3}{3} \right]_0^2 = \frac{8}{3} \]
Multiply the results:
\[ \frac{8}{3} \left( 4 - \frac{5\sqrt{2}}{2} \right) = \frac{32}{3} - \frac{40\sqrt{2}}{6} = \frac{32}{3} - \frac{20\sqrt{2}}{3} \]
\textbf{Final Answers (For Prob 8):}
Integrand box: $4r^2\cos(\theta) - r^2\sin(\theta)$
$A = \pi/4$
$B = \pi/2$
Final Evaluation: $\frac{32}{3} - \frac{20\sqrt{2}}{3}$

\subsection*{Problem 10}
\textbf{Question:} Evaluate $\iint_R \sin(x^2 + y^2) \, dA$, where $R$ is in the first quadrant between the circles $r=2$ and $r=4$.

\textbf{Solution:}
\textbf{Step 1: Set up the bounds.}
First quadrant means $0 \le \theta \le \pi/2$.
The region is between radii 2 and 4, so $2 \le r \le 4$.
\textbf{Step 2: Substitute and integrate.}
\[ \int_0^{\pi/2} \int_2^4 \sin(r^2) \, r \, dr \, d\theta \]
Separate the integrals:
\[ \left( \int_0^{\pi/2} d\theta \right) \left( \int_2^4 r\sin(r^2) \, dr \right) \]
The $\theta$-integral is simply $\pi/2$.
For the $r$-integral, use $u$-substitution. Let $u = r^2$, then $du = 2r \, dr \implies \frac{1}{2}du = r \, dr$.
When $r=2, u=4$. When $r=4, u=16$.
\[ \frac{1}{2} \int_4^{16} \sin(u) \, du = \frac{1}{2} \left[ -\cos(u) \right]_4^{16} = \frac{1}{2} (-\cos(16) + \cos(4)) = \frac{1}{2} (\cos(4) - \cos(16)) \]
Multiply by the $\theta$-integral ($\pi/2$):
\[ \frac{\pi}{2} \times \frac{1}{2} (\cos(4) - \cos(16)) = \frac{\pi}{4}(\cos(4) - \cos(16)) \]
\textbf{Final Answer:} $\frac{\pi}{4}(\cos(4) - \cos(16))$

\subsection*{Problem 11}
\textbf{Question:} Evaluate $\iint_R \frac{y^2}{x^2+y^2} \, dA$, where $R$ is the region between $x^2+y^2=a^2$ and $x^2+y^2=b^2$ with $0 < a < b$.

\textbf{Solution:}
\textbf{Step 1: Determine bounds.}
Since no quadrant is specified, $R$ is a full annulus. $0 \le \theta \le 2\pi$ and $a \le r \le b$.
\textbf{Step 2: Simplify the integrand.}
\[ \frac{y^2}{x^2+y^2} = \frac{r^2\sin^2(\theta)}{r^2} = \sin^2(\theta) \]
Add the $dA = r \, dr \, d\theta$ factor:
\[ \int_0^{2\pi} \int_a^b \sin^2(\theta) \, r \, dr \, d\theta \]
\textbf{Step 3: Integrate.}
Separate the integrals:
\[ \left( \int_0^{2\pi} \sin^2(\theta) \, d\theta \right) \left( \int_a^b r \, dr \right) \]
We know from earlier techniques that integrating $\sin^2(\theta)$ over a full $2\pi$ period yields $\pi$:
\[ \int_0^{2\pi} \frac{1-\cos(2\theta)}{2} \, d\theta = \left[ \frac{\theta}{2} - \frac{\sin(2\theta)}{4} \right]_0^{2\pi} = \pi \]
The $r$-integral is:
\[ \int_a^b r \, dr = \left[ \frac{r^2}{2} \right]_a^b = \frac{b^2 - a^2}{2} \]
Multiply the two results:
\[ \frac{\pi(b^2 - a^2)}{2} \]
\textbf{Final Answer:} $\frac{\pi(b^2 - a^2)}{2}$

\subsection*{Problem 12}
\textbf{Question:} Evaluate $\iint_D e^{-x^2-y^2} \, dA$, bounded by $x = \sqrt{81-y^2}$ and the y-axis.

\textbf{Solution:}
\textbf{Step 1: Understand the region.}
The curve $x = \sqrt{81-y^2}$ represents the right half of a circle of radius $R = 9$ (since $x \ge 0$).
This means the angle spans from $-\pi/2$ to $\pi/2$. The radius goes from $0$ to $9$.
\textbf{Step 2: Set up the integral.}
\[ \int_{-\pi/2}^{\pi/2} \int_0^9 e^{-r^2} \, r \, dr \, d\theta \]
Separate the integrals:
\[ \left( \int_{-\pi/2}^{\pi/2} d\theta \right) \left( \int_0^9 r e^{-r^2} \, dr \right) \]
The $\theta$-integral is $\pi/2 - (-\pi/2) = \pi$.
For the $r$-integral, let $u = -r^2$, so $du = -2r \, dr \implies -\frac{1}{2}du = r \, dr$.
\[ \int_0^9 r e^{-r^2} \, dr = \left[ -\frac{1}{2}e^{-r^2} \right]_0^9 = -\frac{1}{2}e^{-81} - \left(-\frac{1}{2}e^0\right) = \frac{1}{2} - \frac{1}{2}e^{-81} = \frac{1}{2}(1 - e^{-81}) \]
Multiply by the $\theta$-integral result ($\pi$):
\[ \frac{\pi}{2}(1 - e^{-81}) \]
\textbf{Final Answer:} $\frac{\pi}{2}(1 - e^{-81})$

\subsection*{Problem 13}
\textbf{Question:} Evaluate $\iint_D \cos(\sqrt{x^2+y^2}) \, dA$, where $D$ is the disk with center origin and radius 2.

\textbf{Solution:}
\textbf{Step 1: Set up the bounds.}
A full disk means $0 \le \theta \le 2\pi$ and $0 \le r \le 2$.
\textbf{Step 2: Set up the integral.}
\[ \int_0^{2\pi} \int_0^2 \cos(r) \, r \, dr \, d\theta \]
\textbf{Step 3: Integrate using Integration by Parts.}
Separate out the $\theta$-integral, which simply multiplies the result by $2\pi$.
Now, calculate $\int_0^2 r\cos(r) \, dr$.
Let $u = r$, $dv = \cos(r) \, dr$.
Then $du = dr$, $v = \sin(r)$.
Using $\int u \, dv = uv - \int v \, du$:
\[ \int r\cos(r) \, dr = r\sin(r) - \int \sin(r) \, dr = r\sin(r) + \cos(r) \]
Evaluate from 0 to 2:
\[ \left[ 2\sin(2) + \cos(2) \right] - \left[ 0\sin(0) + \cos(0) \right] = 2\sin(2) + \cos(2) - 1 \]
Multiply by the outer integral $2\pi$:
\[ 2\pi(2\sin(2) + \cos(2) - 1) \]
\textbf{Final Answer:} $2\pi(2\sin(2) + \cos(2) - 1)$

\subsection*{Problem 14}
\textbf{Question:} Set up an iterated integral in polar coordinates for the volume of the solid under the surface and above the region $D$. (Limits shown: $0$ to $\pi/2$ and $0$ to $6$). Evaluate it.
\textit{Note: The provided image contains a spinning loading icon, obscuring the exact function graph. Because the explicit function $f(x,y)$ is hidden, I will provide the full framework using the standard function $z = x^2 + y^2$. If your specific problem uses a different function like $z = xy$, simply replace the inner algebraic term.}

\textbf{Solution:}
\textbf{Step 1: Infer bounds.}
The provided structure $\int_0^{\pi/2} \int_0^6$ tells us the region is a quarter-circle in the first quadrant with radius 6. 
\textbf{Step 2: Set up the integral assuming $f(x,y) = x^2 + y^2$.}
In polar coordinates, $x^2 + y^2 = r^2$.
Adding the Jacobian $r$, the integrand is $r^2 \cdot r = r^3$.
\[ \int_0^{\pi/2} \int_0^6 r^3 \, dr \, d\theta \]
\textbf{Step 3: Evaluate.}
\[ \int_0^6 r^3 \, dr = \left[ \frac{r^4}{4} \right]_0^6 = \frac{1296}{4} = 324 \]
Multiply by the angle span:
\[ 324 \times \int_0^{\pi/2} d\theta = 324 \times \frac{\pi}{2} = 162\pi \]
\textbf{Final Answers (assuming $z=x^2+y^2$):}
Integrand box: $r^3$
Final Evaluation: $162\pi$
\textit{(If your function turns out to be $z=xy$, the integrand is $r^3\cos(\theta)\sin(\theta)$ and the result is $162$.)}

\subsection*{Problem 15}
\textbf{Question:} Set up an iterated integral for the volume under $f(x,y) = xy^2$ and above region $D$ (shown as the first-quadrant annulus sector between radii 4 and 5).

\textbf{Solution:}
\textbf{Step 1: Set bounds.}
The graph shows a region in the first quadrant, so $0 \le \theta \le \pi/2$. 
The radius is between 4 and 5, so $4 \le r \le 5$.
\textbf{Step 2: Setup integrand.}
\[ f(x,y) = (r\cos(\theta))(r\sin(\theta))^2 = r^3\cos(\theta)\sin^2(\theta) \]
Multiply by $dA = r \, dr \, d\theta$:
\[ \int_0^{\pi/2} \int_4^5 r^4\cos(\theta)\sin^2(\theta) \, dr \, d\theta \]
\textbf{Step 3: Evaluate.}
Separate the integrals:
\[ \left( \int_0^{\pi/2} \cos(\theta)\sin^2(\theta) \, d\theta \right) \left( \int_4^5 r^4 \, dr \right) \]
For the $\theta$-integral, let $u = \sin(\theta), du = \cos(\theta) \, d\theta$:
\[ \int_0^1 u^2 \, du = \left[ \frac{u^3}{3} \right]_0^1 = \frac{1}{3} \]
For the $r$-integral:
\[ \int_4^5 r^4 \, dr = \left[ \frac{r^5}{5} \right]_4^5 = \frac{3125}{5} - \frac{1024}{5} = \frac{2101}{5} \]
Multiply together:
\[ \frac{1}{3} \times \frac{2101}{5} = \frac{2101}{15} \]
\textbf{Final Answers:}
Integrand box: $r^4\cos(\theta)\sin^2(\theta)$
Final Evaluation: $\frac{2101}{15}$

\subsection*{Problem 16}
\textbf{Question:} Set up an iterated integral for the volume under $f(x,y) = 8x$ and above the region bounded by $r = \sin(\theta)$ in the first quadrant.

\textbf{Solution:}
\textbf{Step 1: Set bounds.}
The bounds are given directly in the problem image: $0 \le \theta \le \pi/2$ and $0 \le r \le \sin(\theta)$.
\textbf{Step 2: Setup integrand.}
\[ 8x \cdot r = 8(r\cos(\theta)) \cdot r = 8r^2\cos(\theta) \]
\[ \int_0^{\pi/2} \int_0^{\sin(\theta)} 8r^2\cos(\theta) \, dr \, d\theta \]
\textbf{Step 3: Evaluate.}
Integrate with respect to $r$ first:
\[ \int_0^{\sin(\theta)} 8r^2\cos(\theta) \, dr = \cos(\theta) \left[ \frac{8r^3}{3} \right]_0^{\sin(\theta)} = \frac{8}{3}\sin^3(\theta)\cos(\theta) \]
Now integrate with respect to $\theta$:
\[ \int_0^{\pi/2} \frac{8}{3}\sin^3(\theta)\cos(\theta) \, d\theta \]
Use substitution $u = \sin(\theta), du = \cos(\theta) \, d\theta$. Limits go from $0$ to $1$.
\[ \frac{8}{3} \int_0^1 u^3 \, du = \frac{8}{3} \left[ \frac{u^4}{4} \right]_0^1 = \frac{8}{3} \left( \frac{1}{4} \right) = \frac{2}{3} \]
\textbf{Final Answers:}
Integrand box: $8r^2\cos(\theta)$
Final Evaluation: $\frac{2}{3}$

\subsection*{Problem 17}
\textbf{Question:} Find the volume under the paraboloid $z = x^2+y^2$ and above the disk $x^2+y^2 \le 9$.

\textbf{Solution:}
\textbf{Step 1: Set bounds.}
The disk covers all angles $0 \le \theta \le 2\pi$ and radius $0 \le r \le 3$.
\textbf{Step 2: Setup integral.}
The function is $z = r^2$. The volume element adds another $r$.
\[ V = \int_0^{2\pi} \int_0^3 (r^2)r \, dr \, d\theta = \int_0^{2\pi} \int_0^3 r^3 \, dr \, d\theta \]
\textbf{Step 3: Evaluate.}
\[ \int_0^3 r^3 \, dr = \left[ \frac{r^4}{4} \right]_0^3 = \frac{81}{4} \]
Multiply by the angular range ($2\pi$):
\[ 2\pi \left( \frac{81}{4} \right) = \frac{81\pi}{2} \]
\textbf{Final Answer:} $\frac{81\pi}{2}$

\subsection*{Problem 18}
\textbf{Question:} Volume inside the sphere $x^2+y^2+z^2=16$ and outside the cylinder $x^2+y^2=1$.

\textbf{Solution:}
\textbf{Step 1: Exploit symmetry.}
The sphere is symmetric above and below the $xy$-plane ($z=0$). We can calculate the volume of the top half and multiply by 2.
The top surface of the sphere is $z = \sqrt{16 - x^2 - y^2} = \sqrt{16 - r^2}$.
\textbf{Step 2: Set bounds.}
The region in the $xy$-plane is a full ring outside radius 1 and inside the sphere's maximum radius (which is 4). 
Bounds: $0 \le \theta \le 2\pi$ and $1 \le r \le 4$.
\textbf{Step 3: Evaluate.}
\[ V_{top} = \int_0^{2\pi} \int_1^4 \sqrt{16-r^2} \, r \, dr \, d\theta \]
Let $u = 16-r^2$, so $du = -2r \, dr \implies -\frac{1}{2}du = r \, dr$.
\[ \int \sqrt{16-r^2} \, r \, dr = -\frac{1}{2} \int u^{1/2} \, du = -\frac{1}{2} \left( \frac{2}{3}u^{3/2} \right) = -\frac{1}{3}(16-r^2)^{3/2} \]
Evaluate from 1 to 4:
\[ \left[ -\frac{1}{3}(16-r^2)^{3/2} \right]_1^4 = -\frac{1}{3}(0 - 15^{3/2}) = \frac{1}{3}(15\sqrt{15}) = 5\sqrt{15} \]
Multiply by $2\pi$ for the angular component: $V_{top} = 10\pi\sqrt{15}$.
Multiply by 2 for the bottom half: $V_{total} = 20\pi\sqrt{15}$.
\textbf{Final Answer:} $20\pi\sqrt{15}$

\subsection*{Problem 19}
\textbf{Question:} Find the volume of a sphere of radius $a$.

\textbf{Solution:}
\textbf{Step 1: Setup upper hemisphere integral.}
Top surface is $z = \sqrt{a^2 - r^2}$. Bounds are $0 \le r \le a$ and $0 \le \theta \le 2\pi$.
\[ V_{top} = \int_0^{2\pi} \int_0^a \sqrt{a^2-r^2} \, r \, dr \, d\theta \]
\textbf{Step 2: Integrate.}
Using substitution $u = a^2-r^2$ similarly to Problem 18:
\[ \int_0^a \sqrt{a^2-r^2} \, r \, dr = \left[ -\frac{1}{3}(a^2-r^2)^{3/2} \right]_0^a = -\frac{1}{3}(0 - a^3) = \frac{a^3}{3} \]
Multiply by $2\pi$ for the angular integration:
\[ V_{top} = \frac{2\pi a^3}{3} \]
Multiply by 2 for the full sphere:
\[ V = \frac{4\pi a^3}{3} \]
\textbf{Final Answer:} $\frac{4\pi a^3}{3}$

\subsection*{Problem 20}
\textbf{Question:} Find the volume bounded by paraboloid $z=3+2x^2+2y^2$ and plane $z=9$ in the first octant.

\textbf{Solution:}
\textbf{Step 1: Find the intersection for radial bounds.}
Set the surfaces equal: $9 = 3 + 2r^2 \implies 6 = 2r^2 \implies r^2 = 3 \implies r = \sqrt{3}$.
Since it's restricted to the first octant, $0 \le \theta \le \pi/2$.
\textbf{Step 2: Setup the integral.}
The height of the solid is the upper surface minus the lower surface: $z = 9 - (3 + 2r^2) = 6 - 2r^2$.
\[ V = \int_0^{\pi/2} \int_0^{\sqrt{3}} (6 - 2r^2) \, r \, dr \, d\theta = \int_0^{\pi/2} \int_0^{\sqrt{3}} (6r - 2r^3) \, dr \, d\theta \]
\textbf{Step 3: Evaluate.}
\[ \int_0^{\sqrt{3}} (6r - 2r^3) \, dr = \left[ 3r^2 - \frac{r^4}{2} \right]_0^{\sqrt{3}} = 3(3) - \frac{9}{2} = 9 - 4.5 = 4.5 = \frac{9}{2} \]
Multiply by the $\theta$-interval ($\pi/2$):
\[ \frac{\pi}{2} \times \frac{9}{2} = \frac{9\pi}{4} \]
\textbf{Final Answer:} $\frac{9\pi}{4}$

\subsection*{Problem 21}
\textbf{Question:} Find the volume above the cone $z = \sqrt{x^2+y^2}$ and below the sphere $x^2+y^2+z^2=49$.

\textbf{Solution:}
\textbf{Step 1: Set up equations in polar.}
Cone: $z = r$.
Sphere: $z = \sqrt{49 - r^2}$.
\textbf{Step 2: Find intersection.}
Set $z = z$: $r = \sqrt{49 - r^2} \implies r^2 = 49 - r^2 \implies 2r^2 = 49 \implies r^2 = 49/2 \implies r = \frac{7}{\sqrt{2}}$.
\textbf{Step 3: Setup integral.}
The volume is the integral of (Sphere - Cone) over the circular region bounded by $r = \frac{7}{\sqrt{2}}$.
\[ V = \int_0^{2\pi} \int_0^{7/\sqrt{2}} \left( \sqrt{49-r^2} - r \right) r \, dr \, d\theta \]
\textbf{Step 4: Evaluate.}
Split into two inner integrals:
1. $\int_0^{7/\sqrt{2}} r\sqrt{49-r^2} \, dr$. Let $u = 49-r^2$, $du = -2r \, dr$. 
Limits map from $49$ down to $49 - \frac{49}{2} = \frac{49}{2}$.
\[ = \left[ -\frac{1}{3}(49-r^2)^{3/2} \right]_0^{7/\sqrt{2}} = -\frac{1}{3}\left(\frac{49}{2}\right)^{3/2} - \left(-\frac{1}{3}(49)^{3/2}\right) \]
\[ = \frac{343}{3} - \frac{343}{3 \times 2\sqrt{2}} = \frac{343}{3} - \frac{343\sqrt{2}}{12} \]
2. $\int_0^{7/\sqrt{2}} r^2 \, dr = \left[ \frac{r^3}{3} \right]_0^{7/\sqrt{2}} = \frac{343}{3 \times 2\sqrt{2}} = \frac{343\sqrt{2}}{12}$
Subtract integral 2 from integral 1:
\[ \left( \frac{343}{3} - \frac{343\sqrt{2}}{12} \right) - \frac{343\sqrt{2}}{12} = \frac{343}{3} - \frac{343\sqrt{2}}{6} \]
Factor out $\frac{343}{3}$:
\[ \frac{343}{3} \left( 1 - \frac{\sqrt{2}}{2} \right) = \frac{343(2-\sqrt{2})}{6} \]
Multiply by $2\pi$:
\[ 2\pi \times \frac{343(2-\sqrt{2})}{6} = \frac{343\pi(2-\sqrt{2})}{3} \]
\textbf{Final Answer:} $\frac{343\pi(2-\sqrt{2})}{3}$

\subsection*{Problem 22}
\textbf{Question:} Find the volume bounded by $z = 10-x^2-y^2$ and $z = 4x^2+4y^2$.

\textbf{Solution:}
\textbf{Step 1: Set up polar functions.}
Top surface: $z = 10 - r^2$.
Bottom surface: $z = 4r^2$.
\textbf{Step 2: Find intersection.}
$10 - r^2 = 4r^2 \implies 5r^2 = 10 \implies r^2 = 2 \implies r = \sqrt{2}$.
\textbf{Step 3: Setup integral.}
\[ V = \int_0^{2\pi} \int_0^{\sqrt{2}} \left( (10 - r^2) - 4r^2 \right) r \, dr \, d\theta = \int_0^{2\pi} \int_0^{\sqrt{2}} (10r - 5r^3) \, dr \, d\theta \]
\textbf{Step 4: Evaluate.}
\[ \int_0^{\sqrt{2}} (10r - 5r^3) \, dr = \left[ 5r^2 - \frac{5r^4}{4} \right]_0^{\sqrt{2}} = 5(2) - \frac{5(4)}{4} = 10 - 5 = 5 \]
Multiply by $2\pi$:
\[ 2\pi(5) = 10\pi \]
\textbf{Final Answer:} $10\pi$

\newpage
\section{Part 3: In-Depth Analysis of Problems and Techniques}

\subsection{A) Problem Types and General Approach}
Categorizing the homework questions reveals three distinct mathematical structures:

\begin{enumerate}
    \item \textbf{Setup Only \& Conversion Problems (e.g., Prob 1-3, 7):}
    These problems test visual recognition. The strategy is to identify the angular sweep (which quadrants does the shape occupy?) to find the limits for $d\theta$, and identify the distance from the origin to find the limits for $dr$.
    
    \item \textbf{Evaluating Definite Integrals with U-Substitution (e.g., Prob 10, 12, 18, 19, 21):}
    Characterized by integrands containing functions composed with $x^2+y^2$ like $e^{-x^2-y^2}$ or $\sin(x^2+y^2)$. 
    \textit{Strategy:} Because the polar conversion adds an $r$ multiplier ($r \, dr \, d\theta$), this perfectly setups up the derivative for a $u$-substitution where $u = r^2$.

    \item \textbf{Volume Between Two Surfaces (e.g., Prob 20, 21, 22):}
    These require finding a 3D volume trapped between a ceiling and a floor. 
    \textit{Strategy:} First, equate the two equations to find the intersection boundary in the xy-plane (this gives the $r$-limits). Second, subtract the bottom function from the top function to create your 2D integrand.
\end{enumerate}

\subsection{B) Key Algebraic and Calculus Manipulations}
\begin{itemize}
    \item \textbf{The "Jacobian" Savior Trick:} The algebraic inclusion of $r$ from $dA = r \, dr \, d\theta$ is not just a requirement; it is a mathematical savior. In Problem 12, evaluating $\iint e^{-x^2-y^2} dx \, dy$ is analytically impossible without advanced techniques. However, changing to polar yields $\int r e^{-r^2} dr$, which is immediately solvable.
    \item \textbf{Trigonometric Power Reduction (Prob 4, 11):} When integrating expressions involving $\sin^2(\theta)$ or $\cos^2(\theta)$, algebraic manipulation via half-angle identities ($\sin^2(\theta) = \frac{1 - \cos(2\theta)}{2}$) is absolutely essential.
    \item \textbf{Separation of Integrals:} Because the bounds of integration in polar circles are strictly constant numbers, and variables can be factored (e.g., $r^2\cos(\theta)$), we can pull the integral apart: $\int \int R(r)\Theta(\theta) \, dr \, d\theta = (\int R(r) \, dr) \times (\int \Theta(\theta) \, d\theta)$. This manipulation was repeatedly used to save time in long calculations.
\end{itemize}

\newpage
\section{Part 4: "Cheatsheet" and Tips for Success}

\begin{tcolorbox}[title=The Polar Coordinates Cheatsheet]
\textbf{Crucial Formulas:}
\begin{itemize}
    \item \textbf{Area Element:} $dA = r \, dr \, d\theta$ (Never forget the $r$!)
    \item \textbf{Conversions:} $x^2 + y^2 = r^2$, \quad $x = r\cos(\theta)$, \quad $y = r\sin(\theta)$
\end{itemize}

\textbf{Shortcuts \& Tricks:}
\begin{itemize}
    \item \textbf{The $2\pi$ Bypass:} If the integrand has no $\theta$ variables and the region is a full circle, you can skip integrating $d\theta$ and simply multiply your $dr$ result by $2\pi$.
    \item \textbf{Identify Quadrant Bounds Instantly:}
        \begin{itemize}
            \item Q1 only: $\theta \in [0, \pi/2]$
            \item Upper Half (Q1 \& Q2): $\theta \in [0, \pi]$
            \item Right Half (Q1 \& Q4): $\theta \in [-\pi/2, \pi/2]$
        \end{itemize}
\end{itemize}

\textbf{Common Pitfalls (Do Not Make These Mistakes):}
\begin{itemize}
    \item \textbf{The Missing "$r$":} Writing $f(r\cos\theta, r\sin\theta) \, dr \, d\theta$ without multiplying the function by $r$.
    \item \textbf{Squaring Errors:} When converting $(x^2+y^2)^2$, remember it equals $(r^2)^2 = r^4$, not just $r^2$.
    \item \textbf{Blindly Using $[0, 2\pi]$:} Always check the prompt for keywords like "first octant" (which restricts bounds to $[0, \pi/2]$) or "first quadrant".
\end{itemize}
\end{tcolorbox}

\newpage
\section{Part 5: Conceptual Synthesis and The "Big Picture"}

\subsection{A) Thematic Connections}
The core theme of this topic is \textbf{Symmetry and Simplification through Transformation}. This maps directly to what we learned previously with $u$-substitution in single-variable calculus. Just as $u$-substitution re-scales the 1D number line to absorb complex inside-functions, Polar Transformation re-maps the 2D grid so that circular boundaries become straight lines in the $(r, \theta)$ plane. We are manipulating the frame of reference so the geometry aligns naturally with our mathematical tools.

\subsection{B) Forward and Backward Links}
\textbf{Backward Links:} This chapter logically extends the fundamental definition of the Riemann sum. We previously summed up infinitesimally thin rectangular pillars ($dx \, dy$) to find volume. We are still doing exactly that, but our "pillars" now have curved faces (wedge shapes). \\
\textbf{Forward Links:} Mastery of the $r \, dr \, d\theta$ conversion is the absolute bedrock for Chapter 15.8: Cylindrical Coordinates ($dz \, r \, dr \, d\theta$) and Chapter 15.9: Spherical Coordinates ($\rho^2 \sin(\phi) \, d\rho \, d\phi \, d\theta$). The intuition you build here on how area stretches ($dA \to r \, dr \, d\theta$) prepares you for calculating the general Jacobian determinant, an essential concept in advanced vector calculus and differential geometry.

\newpage
\section{Part 6: Real-World Application and Modeling}

\subsection{A) Concrete Scenarios}
\begin{enumerate}
    \item \textbf{Quantitative Finance \& Stochastic Calculus:} A quant needs to model the joint probability distribution of two independent, normally distributed assets (e.g., daily returns of a tech stock and a utility stock). The foundational math underlying Brownian Motion and Black-Scholes derivative pricing relies heavily on solving the Gaussian integral over a circular domain to find continuous risk probabilities.
    \item \textbf{Acoustic Engineering:} Designing a parabolic microphone dish. An engineer needs to determine the total energy captured by a circular parabolic dish. The incoming energy density drops off near the edges, modeled as a function of the radius $E(x,y) = I_0 e^{-(x^2+y^2)}$, requiring a polar double integral.
\end{enumerate}

\subsection{B) Model Problem Setup (Finance/Stats)}
Let's model the \textbf{Quantitative Finance} scenario. 
Let $x$ represent the return of Asset A and $y$ represent the return of Asset B. Assume they have a standardized bivariate normal probability density function:
\[ p(x,y) = \frac{1}{2\pi} e^{-\frac{1}{2}(x^2+y^2)} \]
A risk manager wants to know the probability that the combined return profile falls within a severe boundary risk threshold of radius $R = 3$ standard deviations from the mean.
\textbf{Equation to Solve:}
To find this probability, the risk manager must integrate the joint probability density function over the circular domain $D$ defined by $x^2 + y^2 \le 9$.
\[ P = \iint_D \frac{1}{2\pi} e^{-\frac{1}{2}(x^2+y^2)} \, dA \]
Converting to polar to make the integral solvable:
\[ P = \int_0^{2\pi} \int_0^3 \frac{1}{2\pi} e^{-\frac{1}{2}r^2} r \, dr \, d\theta \]
\textit{This integral resolves to precisely $1 - e^{-4.5}$, allowing the firm to perfectly price downside risk.}

\newpage
\section{Part 7: Common Variations and Untested Concepts}

A major application of polar double integrals absent from your homework is calculating \textbf{Center of Mass} and \textbf{Moments of Inertia} for laminas (flat plates) with variable density.

\textbf{Untested Concept: Center of Mass}
If a circular plate $D$ has a density that increases the further you get from the center, given by $\rho(x,y) = \sqrt{x^2+y^2}$, find the mass of the plate if its radius is $R=2$.

\textbf{Brief Explanation:}
Mass is the double integral of density. $m = \iint_D \rho(x,y) \, dA$. In polar, the density function is simply $\rho(r,\theta) = r$. 

\textbf{Worked Example:}
\[ m = \int_0^{2\pi} \int_0^2 (r) \cdot r \, dr \, d\theta \]
\[ m = \int_0^{2\pi} d\theta \int_0^2 r^2 \, dr \]
\[ m = 2\pi \left[ \frac{r^3}{3} \right]_0^2 = 2\pi \left( \frac{8}{3} \right) = \frac{16\pi}{3} \]

\newpage
\section{Part 8: Advanced Diagnostic Testing: "Find the Flaw"}

Below are 5 problems with complete, step-by-step solutions. Each solution contains a single, subtle, fatal flaw. Your job is to act as the professor: find the error, explain it, and solve the problem correctly.

\subsection*{Flawed Problem 1}
\textbf{Question:} Evaluate $\iint_D \cos(x^2+y^2) \, dA$, where $D$ is the region in the first quadrant bounded by $x^2+y^2 = 9$. \\
\textbf{Flawed Solution:} 
1. The bounds are $\theta \in [0, \pi/2]$ and $r \in [0, 3]$. 
2. Substituting $x^2+y^2 = r^2$, the integral becomes: $\int_0^{\pi/2} \int_0^3 \cos(r^2) \, dr \, d\theta$. 
3. This integral requires advanced non-elementary functions (Fresnel integrals), which yields approx $0.621 \times \pi/2$. 

\subsection*{Flawed Problem 2}
\textbf{Question:} Find the area of the region $D$ bounded by the upper half of the circle $x^2+y^2 = 16$. \\
\textbf{Flawed Solution:}
1. The area is $\iint_D 1 \, dA$. 
2. The bounds are $\theta \in [0, \pi]$ and $r \in [0, 4]$.
3. Integral setup: $\int_0^\pi \int_0^4 1 \, dr \, d\theta$. 
4. Inner integral evaluates to $4$. 
5. Outer integral gives $4\pi$. The area is $4\pi$.

\subsection*{Flawed Problem 3}
\textbf{Question:} Evaluate $\iint_D y \, dA$, where $D$ is the right half of the disk $x^2+y^2 \le 4$. \\
\textbf{Flawed Solution:}
1. The region is the right half, so $\theta \in [0, \pi]$. The radius is $r \in [0, 2]$. 
2. Integrand is $y = r\sin(\theta)$. We multiply by $r$ to get $r^2\sin(\theta)$. 
3. $\int_0^\pi \int_0^2 r^2\sin(\theta) \, dr \, d\theta = \int_0^\pi \sin(\theta) \, d\theta \times \int_0^2 r^2 \, dr$. 
4. $[ -\cos(\theta) ]_0^\pi \times [r^3/3]_0^2 = 2 \times (8/3) = 16/3$.

\subsection*{Flawed Problem 4}
\textbf{Question:} Set up the integral for the volume bounded above by $z = x^2+y^2$ and below by the xy-plane, inside $x^2+y^2=25$. \\
\textbf{Flawed Solution:}
1. Paraboloid ceiling: $z = r^2$. Floor: $z=0$. 
2. Intersection bounds: $r = 5$, full circle so $0 \le \theta \le 2\pi$. 
3. Volume integral: $\int_0^{2\pi} \int_0^5 (r^2)^2 \, dr \, d\theta = \int_0^{2\pi} \int_0^5 r^4 \, dr \, d\theta$. 
4. Result: $2\pi [r^5/5]_0^5 = 1250\pi$.

\subsection*{Flawed Problem 5}
\textbf{Question:} Find the volume under $z = 10$ and above $z = 2x^2+2y^2$. \\
\textbf{Flawed Solution:}
1. Find intersection: $10 = 2r^2 \implies r^2 = 5 \implies r = \sqrt{5}$. 
2. Subtract functions to find height: $2r^2 - 10$. 
3. $\int_0^{2\pi} \int_0^{\sqrt{5}} (2r^2 - 10)r \, dr \, d\theta = \int_0^{2\pi} \int_0^{\sqrt{5}} (2r^3 - 10r) \, dr \, d\theta$. 
4. Evaluation leads to $2\pi [ \frac{r^4}{2} - 5r^2 ]_0^{\sqrt{5}} = 2\pi [ 12.5 - 25 ] = -25\pi$. Volume cannot be negative, take the absolute value: $25\pi$.

\vspace{1cm}
\begin{center}
\rule{0.5\linewidth}{0.5pt}
\end{center}
\vspace{1cm}

\subsection*{Answers to "Find the Flaw"}
\textbf{Flaw 1:} The student forgot to multiply by the Jacobian $r$ from $dA = r \, dr \, d\theta$. 
\textit{Correct:} The integral should be $\int_0^{\pi/2} \int_0^3 \cos(r^2) \, r \, dr \, d\theta$, which easily resolves via u-substitution to $\frac{\pi}{4}\sin(9)$.

\textbf{Flaw 2:} The student forgot the extra $r$ when setting up the integral for area ($dA = r \, dr \, d\theta$). 
\textit{Correct:} The inner integral should be $\int_0^4 r \, dr = 8$. Multiplied by $\pi$, the true area is $8\pi$.

\textbf{Flaw 3:} The student used incorrect bounds for the "right half" of a disk. $[0, \pi]$ represents the top half, not the right half. 
\textit{Correct:} The bounds for the right half are $[-\pi/2, \pi/2]$. Integrating $\sin(\theta)$ over these bounds yields $0$.

\textbf{Flaw 4:} The student squared the function $r^2$ unnecessarily. The surface is $z = r^2$. We multiply by $r$ (from $dA$), resulting in $r^3$, not $(r^2)^2 = r^4$. 
\textit{Correct:} $\int_0^{2\pi} \int_0^5 r^3 \, dr \, d\theta = 2\pi [r^4/4]_0^5 = \frac{625\pi}{2}$.

\textbf{Flaw 5:} The student subtracted the top function from the bottom function ($Floor - Ceiling$), resulting in a negative volume. You cannot just "take the absolute value" as a math fix, it indicates improper structural setup. 
\textit{Correct:} The height must be $10 - 2r^2$. The correct integral evaluates naturally to $+25\pi$.

\end{document}